\question[4]
Viele Apps auf dem Smartphone, z.\,B.\ Wetter-Apps, Navigations-Apps 
oder soziale Netzwerke, fragen bei der Installation nach der 
Berechtigung, auf den Standort (Geodaten) des Geräts zuzugreifen.

\begin{enumerate}
    \item Erkläre, wie Smartphones den eigenen Standort technisch 
    bestimmen können.
    \item Nenne zwei legitime Gründe, warum eine App Geodaten 
    benötigen könnte.
    \item Beschreibe ein Risiko, das entsteht, wenn Geodaten über 
    einen längeren Zeitraum gesammelt und gespeichert werden.
\end{enumerate}

\begin{solutionbox}{8cm}

\textbf{1.} Smartphones bestimmen ihren Standort meist über GPS 
(Satellitensignale), können ihn aber auch über die Position naher 
WLAN-Netze, die Entfernung zu Mobilfunkmasten oder eine Kombination 
dieser Verfahren ermitteln, wenn kein GPS-Signal verfügbar ist.

\textbf{2.} Beispiele: Eine Navigations-App benötigt den Standort, um 
eine Route berechnen zu können. Eine Wetter-App benötigt ihn, um das 
Wetter für den aktuellen Aufenthaltsort anzuzeigen. (Weitere gültige 
Beispiele: Standortmarkierung bei Fotos, Finden von Geschäften in der 
Nähe.)

\textbf{3.} Werden Geodaten über längere Zeit gesammelt, lässt sich 
aus den Bewegungsmustern ein detailliertes Profil erstellen, aus dem 
sich z.\,B.\ Wohnort, Arbeitsplatz, Tagesablauf oder besuchte Orte 
(etwa Arztpraxen) ablesen lassen. Gelangen diese Daten in falsche 
Hände oder werden sie mit anderen Daten verknüpft, können daraus 
sehr persönliche Rückschlüsse gezogen und die Betroffenen z.\,B.\ 
überwacht, erpresst oder gezielt beeinflusst werden (z.\,B.\ durch 
personalisierte Werbung).

\end{solutionbox}